\documentclass[../../../../master.tex]{subfiles}
\begin{document}
\begingroup\sloppy\emergencystretch=3em

\monographleaf{Domain 06: Reinforcement Learning \& Alignment}{Bandits, Markov Decision Processes \& Value-Based RL}{Monte Carlo, Temporal-Difference, DQN \& Function Approximation}{Monte Carlo, SARSA, Q-Learning, DQN, Eligibility Traces \& the Deadly Triad}{Sutton \& Barto Ch.~5--7, 9--12; Murphy RL Sec.~2.3--2.5; Murphy Adv. Sec.~35.2; Prince Sec.~19.4}

\section{Problem Formulation \& Unified Notation}
\label{sec:td:formulation}

In model-free reinforcement learning (RL), an agent interacts with an unknown Markov Decision Process (MDP) $\mathcal{M} = (\States, \Actions, \Trans, \Rewards, \gamma)$ through sampled trajectories $(S_0, A_0, R_1, S_1, A_1, R_2, \dots, S_T)$ without access to the transition kernel $\Trans(s' \mid s, a)$ or expected reward function \bookcite{reinforcement_learning_introduction_richard_s_su}{Ch.~5, PDF p.~113} \bookcite{reinforement_learning_kevin_murphy}{Sec.~2.3, p.~35}. Given a policy $\pi(a \mid s)$, the agent seeks to estimate the state-value function $v_\pi(s) = \E_\pi[G_t \mid S_t = s]$ or action-value function $q_\pi(s, a) = \E_\pi[G_t \mid S_t = s, A_t = a]$ (prediction), and ultimately discover an optimal policy $\pi^*$ maximizing the expected discounted return (control):
\begin{equation}
  G_t \;\doteq\; \sum_{k=0}^{T - t - 1} \gamma^k R_{t+k+1}, \qquad \gamma \in [0, 1],
  \label{eq:td:mc_return}
\end{equation}
where $T = \infty$ for continuing tasks ($\gamma < 1$) and $T < \infty$ for episodic tasks. Model-free value estimation spans a fundamental bias--variance--memory spectrum between pure Monte Carlo (MC) rollouts, one-step Temporal-Difference ($\mathrm{TD}(0)$) bootstrapping, $n$-step truncated returns, and $\mathrm{TD}(\lambda)$ eligibility traces. When tabular lookup is replaced by parametric function approximation $\hat{v}(s, \vecb{w})$ or $Q_{\vecb{w}}(s, a)$ under off-policy sampling, the interaction of approximation, bootstrapping, and off-policy distribution shift forms the \emph{Deadly Triad}, necessitating architectural and algorithmic stabilizers such as Gradient-TD, Experience Replay, Target Networks, and Double/Dueling DQN \bookcite{reinforcement_learning_introduction_richard_s_su}{Ch.~11, PDF p.~286} \bookcite{reinforement_learning_kevin_murphy}{Sec.~2.5.2, p.~42}.

\begin{crossbook}[Cross-Book Notation \& Pedagogical Reconciliation]
Source texts adopt slightly different indexing conventions for rewards, step sizes, and network parameters. We standardize on Sutton \& Barto's $R_{t+1}$ transition indexing while mapping Murphy's and Prince's conventions explicitly in Table~\ref{tab:td:notation}.
\vspace{0.4em}
\begin{center}
\small
\begin{tabularx}{\textwidth}{@{}l l l l X@{}}
\toprule
\textbf{Concept} & \textbf{Sutton \& Barto} & \textbf{Murphy (RL / Adv.)} & \textbf{Prince} & \textbf{Unified Monograph Convention} \\
\midrule
Step-size / LR & $\alpha$ & $\eta$ & $\alpha$ & $\alpha \in (0, 1]$ (TD/MC), $\eta$ for SGD \\
Reward at $t \to t+1$ & $R_{t+1}$ & $r_t$ or $r_{t+1}$ & $r[s_t, a_t]$ & $R_{t+1}$ (emitted upon $S_t, A_t \to S_{t+1}$) \\
Value parameters & $\vecb{w} \in \R^d$ & $\vecb{w}$ & $\vecb{\phi}$ & $\vecb{w}$ (linear/general), $\vecb{\theta}$ (deep Q-net) \\
Target net params & --- & $\vecb{w}^-$ & $\vecb{\phi}^-$ & $\vecb{\theta}^-$ (frozen / Polyak-averaged) \\
Action-value fn. & $\hat{q}(s, a, \vecb{w})$ & $Q_{\vecb{w}}(s, a)$ & $q[s_t, a_t, \vecb{\phi}]$ & $Q_{\vecb{\theta}}(s, a)$ (or $q_\pi(s, a)$ tabular) \\
State distribution & $\mu(s)$ (on/off-policy) & $d^\pi(s)$ / $\mathcal{D}$ & $\mathcal{D}$ (replay) & $\mu(s)$ for state weighting, $\mathcal{D}$ for buffer \\
\bottomrule
\end{tabularx}
\end{center}
\label{tab:td:notation}
\end{crossbook}

\section{Core Mathematical Derivations \& Theorems}
\label{sec:td:derivations}

\subsection{Monte Carlo Prediction, Importance Sampling \& TD(0) Bootstrapping}

Monte Carlo methods estimate $v_\pi(s)$ from complete episodic returns $G_t$, avoiding any reliance on transition dynamics \bookcite{reinforcement_learning_introduction_richard_s_su}{Ch.~5, PDF p.~114}. In \emph{first-visit MC}, only the first occurrence of state $s$ in an episode contributes to the sample average; in \emph{every-visit MC}, all visits contribute. Both are unbiased in the limit (first-visit is unbiased per sample: $\E[G_t \mid S_t = s] = v_\pi(s)$), but suffer from high variance because $G_t$ accumulates stochasticity across the entire remaining trajectory \bookcite{probabilisticml_kevinmurphy_advancedtopics}{Sec.~35.2.1, PDF p.~1184}.

When evaluating a target policy $\pi$ using episodes generated by a distinct exploratory behavior policy $b$ satisfying the \emph{coverage assumption} ($\pi(a \mid s) > 0 \implies b(a \mid s) > 0$), we weight each return by the trajectory importance-sampling (IS) ratio \bookcite{reinforcement_learning_introduction_richard_s_su}{Ch.~5, PDF pp.~125--127}:
\begin{equation}
  \rho_{t:T-1} \;\doteq\; \prod_{k=t}^{T-1} \frac{\pi(A_k \mid S_k)}{b(A_k \mid S_k)}.
  \label{eq:td:is_ratio}
\end{equation}
\emph{Ordinary IS} averages $\rho_{t:T-1} G_t$ over visits $\mathcal{J}(s)$: $V(s) = \frac{\sum_{t \in \mathcal{J}(s)} \rho_{t:T-1} G_t}{\abs{\mathcal{J}(s)}}$, which is unbiased ($\E_b[\rho_{t:T-1} G_t \mid S_t = s] = v_\pi(s)$) but can have infinite variance when $\rho_{t:T-1}$ is unbounded \bookcite{reinforcement_learning_introduction_richard_s_su}{Ch.~5, PDF pp.~128--129}. \emph{Weighted IS} normalizes by the sum of ratios:
\begin{equation}
  V(s) \;\doteq\; \frac{\sum_{t \in \mathcal{J}(s)} \rho_{t:T-1} G_t}{\sum_{t \in \mathcal{J}(s)} \rho_{t:T-1}},
\end{equation}
introducing finite-sample bias that vanishes asymptotically ($O(1/N)$) while drastically reducing variance.

In contrast, \emph{Temporal-Difference} ($\mathrm{TD}(0)$) learning updates $V(S_t)$ immediately after a single transition $(S_t, A_t, R_{t+1}, S_{t+1})$ by bootstrapping from the current estimate $V(S_{t+1})$ \bookcite{reinforcement_learning_introduction_richard_s_su}{Ch.~6, PDF p.~141} \bookcite{reinforement_learning_kevin_murphy}{Sec.~2.3.2, p.~36}:
\begin{equation}
  V(S_t) \;\leftarrow\; V(S_t) + \alpha \underbrace{\bigl[ R_{t+1} + \gamma V(S_{t+1}) - V(S_t) \bigr]}_{\delta_t \;\text{(TD error)}}.
  \label{eq:td:td0}
\end{equation}

\monographfig[0.86\textwidth]{reinforement_learning_kevin_murphy_p37_fig1.png}{Backup diagrams contrasting sample-based full-depth Monte Carlo (left), sample-based one-step Temporal-Difference bootstrapping (center), and full-width one-step Dynamic Programming expectation (right) \bookcite{reinforement_learning_kevin_murphy}{Fig.~2.3, p.~37}.}{fig:td:backup_spectrum}

\begin{intuition}[Why TD(0) Outperforms MC Under Batch Training]
In finite-sample batch training, MC minimizes the empirical squared error on observed returns $\sum_t (G_t - V(S_t))^2$, whereas $\mathrm{TD}(0)$ converges to the \emph{certainty-equivalence} estimate: the exact value function of the maximum-likelihood MDP inferred from the observed transitions \bookcite{reinforcement_learning_introduction_richard_s_su}{Ch.~6, PDF pp.~148--150}. By exploiting the Markov property rather than ignoring state-transition structure, $\mathrm{TD}(0)$ achieves strictly lower asymptotic estimation error on Markovian environments despite its transient bootstrapping bias.
\end{intuition}

\subsection{TD Control: SARSA, Q-Learning, Expected Sarsa \& Double Learning}

Extending $\mathrm{TD}(0)$ to action values $Q(s, a)$ yields three canonical tabular control algorithms, distinguished by how they construct the bootstrapping target \bookcite{reinforcement_learning_introduction_richard_s_su}{Ch.~6, PDF pp.~151--156}:
\begin{itemize}
  \item \textbf{SARSA (On-Policy):} Samples $A_{t+1} \sim \pi(\cdot \mid S_{t+1})$ from the same $\epsilon$-greedy policy used for behavior:
  \begin{equation}
    Q(S_t, A_t) \;\leftarrow\; Q(S_t, A_t) + \alpha \bigl[ R_{t+1} + \gamma Q(S_{t+1}, A_{t+1}) - Q(S_t, A_t) \bigr].
    \label{eq:td:sarsa}
  \end{equation}
  \item \textbf{Q-Learning (Off-Policy):} Directly approximates $q_*$ independently of the exploratory behavior policy by bootstrapping from the greedy maximum:
  \begin{equation}
    Q(S_t, A_t) \;\leftarrow\; Q(S_t, A_t) + \alpha \bigl[ R_{t+1} + \gamma \max_{a'} Q(S_{t+1}, a') - Q(S_t, A_t) \bigr].
    \label{eq:td:qlearning}
  \end{equation}
  \item \textbf{Expected Sarsa:} Replaces the single sampled next action $A_{t+1}$ with its expectation under target policy $\pi$, eliminating sampling variance over $A_{t+1}$:
  \begin{equation}
    Q(S_t, A_t) \;\leftarrow\; Q(S_t, A_t) + \alpha \Bigl[ R_{t+1} + \gamma \sum_{a'} \pi(a' \mid S_{t+1}) Q(S_{t+1}, a') - Q(S_t, A_t) \Bigr].
    \label{eq:td:expected_sarsa}
  \end{equation}
\end{itemize}
Notice that Q-learning is a special case of Expected Sarsa when $\pi$ is the deterministic greedy policy with respect to $Q$.

\begin{coremath}[Maximization Bias \& Double Q-Learning]
In Q-learning, the target $\max_{a'} Q(S_{t+1}, a')$ uses the same noisy estimates both to select the maximizing action and to evaluate it. By Jensen's inequality for the convex $\max$ operator on random variables $\{X_a\}_{a \in \Actions}$ with $\E[X_a] = q(a)$:
\begin{equation}
  \E\bigl[\max_{a} X_a\bigr] \;\ge\; \max_{a} \E[X_a] \;=\; \max_{a} q(a),
\end{equation}
with strict inequality whenever multiple actions have noisy estimates near the maximum \bookcite{reinforcement_learning_introduction_richard_s_su}{Ch.~6, PDF pp.~156--157} \bookcite{understandingdl_simonprince}{Sec.~19.4.2, p.~401}. This positive \emph{maximization bias} causes Q-learning to overestimate risky/high-variance suboptimal actions.

\textbf{Double Q-Learning} decouples action selection from evaluation by maintaining two independent estimators $Q_1, Q_2$ \bookcite{reinforcement_learning_introduction_richard_s_su}{Ch.~6, PDF pp.~157--158} \bookcite{probabilisticml_kevinmurphy_advancedtopics}{Sec.~35.2.5, PDF p.~1188}. At each step, with probability $0.5$ we update $Q_1$ using $Q_2$ to evaluate $Q_1$'s greedy action (and symmetrically for $Q_2$):
\begin{equation}
  Q_1(S_t, A_t) \;\leftarrow\; Q_1(S_t, A_t) + \alpha \Bigl[ R_{t+1} + \gamma Q_2\bigl(S_{t+1}, \argmax_{a'} Q_1(S_{t+1}, a')\bigr) - Q_1(S_t, A_t) \Bigr].
  \label{eq:td:double_q}
\end{equation}
Since $Q_2$ is independent of the noise that selected $a^* = \argmax_{a'} Q_1(S_{t+1}, a')$, $\E[Q_2(S_{t+1}, a^*)] = q(S_{t+1}, a^*)$ is unbiased for that action's true value.
\end{coremath}

\monographfig[0.74\textwidth]{reinforcement_learning_introduction_richard_s_su_p157_fig1.png}{Empirical demonstration of maximization bias on a two-state episodic MDP: standard Q-learning heavily over-selects the suboptimal \texttt{left} action entering state $B$ (whose actions have $\Normal(-0.1, 1)$ rewards), whereas Double Q-learning eliminates the bias \bookcite{reinforcement_learning_introduction_richard_s_su}{Fig.~6.5, PDF p.~157}.}{fig:td:max_bias}

\subsection{n-Step Bootstrapping \& TD($\lambda$) Eligibility Traces}

Between $\mathrm{TD}(0)$ ($1$-step) and MC ($\infty$-step) lies the $n$-step return \bookcite{reinforcement_learning_introduction_richard_s_su}{Ch.~7, PDF p.~165}:
\begin{equation}
  G_{t:t+n} \;\doteq\; R_{t+1} + \gamma R_{t+2} + \dots + \gamma^{n-1} R_{t+n} + \gamma^n V_{t+n-1}(S_{t+n}), \qquad n \ge 1,
  \label{eq:td:nstep_return}
\end{equation}
with $G_{t:t+n} \doteq G_t$ if $t + n \ge T$. Rather than committing to a single horizon $n$, the \emph{forward-view $\lambda$-return} forms an exponentially decaying convex combination of all $n$-step returns with trace-decay parameter $\lambda \in [0, 1]$ \bookcite{reinforcement_learning_introduction_richard_s_su}{Ch.~12, PDF pp.~310--311} \bookcite{reinforement_learning_kevin_murphy}{Sec.~2.3.3, p.~36}:
\begin{equation}
  G_t^\lambda \;\doteq\; (1 - \lambda) \sum_{n=1}^{\infty} \lambda^{n-1} G_{t:t+n} \;=\; (1 - \lambda)\sum_{n=1}^{T-t-1} \lambda^{n-1} G_{t:t+n} + \lambda^{T-t-1} G_t.
  \label{eq:td:lambda_return}
\end{equation}
Setting $\lambda = 0$ recovers $G_t^0 = G_{t:t+1}$ ($\mathrm{TD}(0)$), while $\lambda = 1$ recovers $G_t^1 = G_t$ (MC).

\begin{coremath}[Forward--Backward Equivalence of $\mathrm{TD}(\lambda)$]
Computing $G_t^\lambda$ directly requires waiting until episode termination $T$ (non-causal). \textbf{Backward-view $\mathrm{TD}(\lambda)$} achieves causal, step-by-step updates using an \emph{eligibility trace} vector $\vecb{z}_t \in \R^d$ that tracks recently active feature gradients \bookcite{reinforcement_learning_introduction_richard_s_su}{Ch.~12, PDF pp.~314--315}:
\begin{equation}
  \vecb{z}_{-1} = \vecb{0}, \qquad \vecb{z}_t \;\doteq\; \gamma \lambda \vecb{z}_{t-1} + \nabla_{\vecb{w}} \hat{v}(S_t, \vecb{w}_t),
  \label{eq:td:eligibility_trace}
\end{equation}
paired with the semi-gradient update driven by the one-step TD error $\delta_t = R_{t+1} + \gamma \hat{v}(S_{t+1}, \vecb{w}_t) - \hat{v}(S_t, \vecb{w}_t)$:
\begin{equation}
  \vecb{w}_{t+1} \;\doteq\; \vecb{w}_t + \alpha \delta_t \vecb{z}_t.
  \label{eq:td:td_lambda_update}
\end{equation}
For linear function approximation $\hat{v}(s, \vecb{w}) = \vecb{w}\T \vecb{x}(s)$ under offline (end-of-episode) weight accumulation, the sum of backward-view updates equals the sum of forward-view $\lambda$-return updates:
\begin{equation}
  \sum_{t=0}^{T-1} \delta_t \vecb{z}_t \;=\; \sum_{t=0}^{T-1} \bigl(G_t^\lambda - \vecb{w}\T \vecb{x}(S_t)\bigr) \vecb{x}(S_t),
\end{equation}
verified by telescoping the geometric series $G_t^\lambda - \hat{v}(S_t) = \sum_{k=t}^{T-1} (\gamma \lambda)^{k-t} \delta_k$ and exchanging summation indices \bookcite{reinforcement_learning_introduction_richard_s_su}{Ch.~12, PDF pp.~315--316}. For online step-by-step updates, \emph{True Online $\mathrm{TD}(\lambda)$} eliminates the discrepancy via Dutch traces \bookcite{reinforcement_learning_introduction_richard_s_su}{Ch.~12, PDF p.~321}.
\end{coremath}

\subsection{Function Approximation, Semi-Gradient TD \& The Deadly Triad}

When $\abs{\States}$ is large or continuous, we parameterize $\hat{v}(s, \vecb{w})$ and minimize the weighted Value Error $\overline{\mathrm{VE}}(\vecb{w}) \doteq \norm{\hat{v}(\cdot, \vecb{w}) - v_\pi}_\mu^2 = \sum_{s \in \States} \mu(s) [\hat{v}(s, \vecb{w}) - v_\pi(s)]^2$ \bookcite{reinforcement_learning_introduction_richard_s_su}{Ch.~9, PDF pp.~221--222}. Because $v_\pi(S_t)$ is unknown, \emph{semi-gradient $\mathrm{TD}(0)$} substitutes the bootstrapping target $U_t = R_{t+1} + \gamma \hat{v}(S_{t+1}, \vecb{w}_t)$ and differentiates only the prediction $\hat{v}(S_t, \vecb{w}_t)$, ignoring the dependence of $U_t$ on $\vecb{w}_t$:
\begin{equation}
  \vecb{w}_{t+1} \;\doteq\; \vecb{w}_t + \alpha \bigl[ R_{t+1} + \gamma \hat{v}(S_{t+1}, \vecb{w}_t) - \hat{v}(S_t, \vecb{w}_t) \bigr] \nabla_{\vecb{w}} \hat{v}(S_t, \vecb{w}_t).
  \label{eq:td:semi_grad_td}
\end{equation}
Under \emph{on-policy} sampling ($\mu = d^\pi$) and linear features $\hat{v}(s, \vecb{w}) = \vecb{w}\T \vecb{x}(s)$, semi-gradient TD(0) converges reliably to the TD fixed point $\vecb{w}_{\mathrm{TD}}$ where the \emph{Projected Bellman Error} vanishes: $\mathrm{PBE}(\vecb{w}) \doteq \norm{\Pi \overline{\delta}_{\vecb{w}}}_\mu^2 = 0$, with $\Pi = \mat{X}(\mat{X}\T \mat{D} \mat{X})^{-1} \mat{X}\T \mat{D}$ and $\mat{D} = \diag(\mu)$ \bookcite{reinforcement_learning_introduction_richard_s_su}{Ch.~11, PDF pp.~288--290}.

However, when \emph{off-policy} sampling ($\mu \ne d^\pi$) is combined with bootstrapping and function approximation, the semi-gradient operator $\Pi \Bellman_\pi$ is no longer guaranteed to be a contraction in $\norm{\cdot}_\mu$, causing weights to diverge to infinity---even in linear MDPs such as \textbf{Baird's 7-state counterexample} \bookcite{reinforcement_learning_introduction_richard_s_su}{Ch.~11, PDF pp.~282--287} \bookcite{reinforement_learning_kevin_murphy}{Sec.~2.5.2.5, p.~43}.

\monographfig[0.86\textwidth]{reinforement_learning_kevin_murphy_p43_fig1.png}{\textbf{Baird's Counterexample:} (a) A 7-state MDP with target policy $\pi(\text{solid} \mid \cdot) = 1$, behavior policy $b(\text{dashed} \mid \cdot) = 6/7$, zero rewards, and $\gamma = 0.99$. Although $\vecb{w} = \vecb{0}$ represents the exact $v_\pi \equiv 0$, (b) off-policy semi-gradient $\mathrm{TD}(0)$ diverges exponentially \bookcite{reinforement_learning_kevin_murphy}{Fig.~2.6, p.~43}.}{fig:td:baird_divergence}

\begin{coremath}[Gradient-TD (GTD2 / TDC) Stabilization]
To obtain $O(d)$ true stochastic gradient descent on the Projected Bellman Error under arbitrary off-policy distributions, \textbf{Gradient-TD} rewrites $\nabla_{\vecb{w}} \mathrm{PBE}(\vecb{w})$ using expectations under behavior policy $b$ \bookcite{reinforcement_learning_introduction_richard_s_su}{Ch.~11, PDF pp.~300--302}:
\begin{equation}
  \nabla_{\vecb{w}} \mathrm{PBE}(\vecb{w}) \;=\; 2 \E\bigl[\rho_t (\gamma \vecb{x}_{t+1} - \vecb{x}_t)\vecb{x}_t\T\bigr] \bigl(\E[\vecb{x}_t \vecb{x}_t\T]\bigr)^{-1} \E[\rho_t \delta_t \vecb{x}_t].
  \label{eq:td:pbe}
\end{equation}
Because the first and third expectations both depend on $\vecb{x}_{t+1}$, multiplying single-step samples is biased. GTD2/TDC circumvents this double-sampling issue by maintaining an auxiliary $d$-vector $\vecb{v}_t \approx (\E[\vecb{x}_t \vecb{x}_t\T])^{-1} \E[\rho_t \delta_t \vecb{x}_t]$ via a fast-timescale LMS update ($\vecb{v}_{t+1} = \vecb{v}_t + \beta \rho_t (\delta_t - \vecb{v}_t\T \vecb{x}_t) \vecb{x}_t$) and updating $\vecb{w}$ on a slower timescale:
\begin{equation}
  \vecb{w}_{t+1} \;=\; \vecb{w}_t + \alpha \rho_t \bigl(\delta_t \vecb{x}_t - \gamma \vecb{x}_{t+1} (\vecb{x}_t\T \vecb{v}_t)\bigr) \qquad (\text{TDC}).
  \label{eq:td:gtd2}
\end{equation}
\end{coremath}

\section{Algorithmic Implementation \& Comparative Synthesis}
\label{sec:td:algorithms}

Scaling Q-learning to high-dimensional observations (e.g., Atari $84 \times 84 \times 4$ frame stacks) via deep neural networks $Q_{\vecb{\theta}}(s, a)$---termed \emph{Fitted Q-Iteration} or \emph{Deep Q-Networks (DQN)}---triggers severe Deadly Triad oscillations because adjacent frames are strongly correlated and the regression target $R_{t+1} + \gamma \max_{a'} Q_{\vecb{\theta}}(S_{t+1}, a')$ moves on every gradient step \bookcite{understandingdl_simonprince}{Sec.~19.4, pp.~399--401} \bookcite{probabilisticml_kevinmurphy_advancedtopics}{Sec.~35.2.6, PDF p.~1188}. Modern deep value-based RL stabilizes training through six synergistic mechanisms unified in \textbf{Rainbow} \bookcite{reinforement_learning_kevin_murphy}{Sec.~2.5.4, pp.~45--48}:
\begin{enumerate}
  \item \textbf{Experience Replay ($\mathcal{D}$):} Stores $(s, a, r, s', \text{done})$ tuples and samples minibatches uniformly or via \textbf{Prioritized Experience Replay (PER)}, where transition $i$ has priority $p_i = (\abs{\delta_i} + \epsilon_{\text{per}})^\alpha$ and importance-sampling correction weights $w_i = (N \cdot P(i))^{-\beta}$ \bookcite{reinforement_learning_kevin_murphy}{Sec.~2.5.2.3--2.5.2.4, p.~42}.
  \item \textbf{Frozen Target Network ($\vecb{\theta}^-$):} Decouples target evaluation from online parameter updates by holding $\vecb{\theta}^-$ fixed for $C$ steps (or Polyak averaging $\vecb{\theta}^- \leftarrow \tau \vecb{\theta} + (1 - \tau)\vecb{\theta}^-$), yielding the DQN loss \bookcite{probabilisticml_kevinmurphy_advancedtopics}{Eq.~(35.14), PDF p.~1189}:
  \begin{equation}
    \mathcal{L}_{\mathrm{DQN}}(\vecb{\theta}) \;=\; \E_{(s, a, r, s') \sim \mathcal{D}} \Bigl[\bigl(r + \gamma \max_{a'} Q_{\vecb{\theta}^-}(s', a') - Q_{\vecb{\theta}}(s, a)\bigr)^2\Bigr].
    \label{eq:td:dqn_loss}
  \end{equation}
  \item \textbf{Double DQN Target:} Eliminates maximization bias by selecting the greedy action with $\vecb{\theta}$ and evaluating it with $\vecb{\theta}^-$ \bookcite{reinforement_learning_kevin_murphy}{Sec.~2.5.3.2, p.~45}:
  \begin{equation}
    y_i^{\mathrm{DDQN}} \;=\; r_i + \gamma (1 - \text{done}_i)\, Q_{\vecb{\theta}^-}\!\bigl(s'_i, \argmax_{a'} Q_{\vecb{\theta}}(s'_i, a')\bigr).
    \label{eq:td:double_dqn_target}
  \end{equation}
  \item \textbf{Dueling Architecture:} Factorizes $Q_{\vecb{\theta}}(s, a)$ into a state-value stream $V_{\vecb{\theta}}(s)$ and an advantage stream $A_{\vecb{\theta}}(s, a)$, enforcing identifiability ($\E_{\pi^*}[A^*] = 0$) via mean subtraction \bookcite{reinforement_learning_kevin_murphy}{Sec.~2.5.4.2, pp.~46--47}:
  \begin{equation}
    Q_{\vecb{\theta}}(s, a) \;=\; V_{\vecb{\theta}}(s) + \Bigl(A_{\vecb{\theta}}(s, a) - \frac{1}{\abs{\Actions}} \sum_{a'} A_{\vecb{\theta}}(s, a')\Bigr).
    \label{eq:td:dueling_q}
  \end{equation}
  \item \textbf{Multi-step ($n=3$) Returns, Noisy Nets \& Distributional RL (C51):} Reduce bootstrapping bias, provide state-dependent exploration, and learn the full return distribution rather than only its expectation \bookcite{reinforement_learning_kevin_murphy}{Sec.~2.5.4.6, p.~48}.
\end{enumerate}

\begin{algorithm}[htbp]
\caption{Double DQN with Prioritized Experience Replay \& Target Network}
\label{alg:td:double_dqn}
\KwIn{Environment $\mathcal{M}$, replay buffer $\mathcal{D}$ (capacity $N$), online weights $\vecb{\theta}$, target weights $\vecb{\theta}^- \leftarrow \vecb{\theta}$, step size $\eta$, discount $\gamma$, target sync period $C$, PER exponents $(\alpha, \beta)$.}
\KwOut{Trained action-value network $Q_{\vecb{\theta}}(s, a)$.}
\For{episode $= 1, \dots, E$}{
  Observe initial state $s_0$\;
  \For{$t = 0, 1, \dots$ until terminal}{
    Select $a_t$ via $\epsilon$-greedy policy in $Q_{\vecb{\theta}}(s_t, \cdot)$\;
    Execute $a_t$, observe $r_{t+1}, s_{t+1}, \text{done}_t$, and store $(s_t, a_t, r_{t+1}, s_{t+1}, \text{done}_t)$ in $\mathcal{D}$ with $p_t = \max_{j} p_j$\;
    Sample minibatch $\{i_k\}_{k=1}^B \sim P(i) = p_i^\alpha / \sum_j p_j^\alpha$\;
    \For{each sample $i$ in minibatch}{
      Compute Double DQN target $y_i \leftarrow r_i + \gamma(1 - \text{done}_i) Q_{\vecb{\theta}^-}(s'_i, \argmax_{a'} Q_{\vecb{\theta}}(s'_i, a'))$\;
      Compute TD error $\delta_i \leftarrow y_i - Q_{\vecb{\theta}}(s_i, a_i)$ and IS weight $w_i \leftarrow (N \cdot P(i))^{-\beta}$\;
      Update priority $p_i \leftarrow (\abs{\delta_i} + \epsilon_{\text{per}})^\alpha$\;
    }
    Update $\vecb{\theta} \leftarrow \vecb{\theta} - \eta \frac{1}{B} \sum_{i} w_i \, \nabla_{\vecb{\theta}} \bigl[\tfrac{1}{2}\delta_i^2\bigr]$ (normalized by $\max_j w_j$)\;
    \lIf{$t \bmod C = 0$}{$\vecb{\theta}^- \leftarrow \vecb{\theta}$}
  }
}
\end{algorithm}

\begin{crossbook}[Cross-Book Synthesis: Sample Efficiency, Goal Relabeling \& Modern DQN]
While Sutton \& Barto emphasize theoretical convergence guarantees (e.g., Gradient-TD and Emphatic-TD) for linear off-policy prediction \bookcite{reinforcement_learning_introduction_richard_s_su}{Ch.~11, PDF pp.~300--304}, Murphy RL bridges classical control to modern sample-efficient deep RL: \emph{Bigger, Better, Faster (BBF)} achieves state-of-the-art Atari-100k performance by combining high update-to-data (UTD) ratios, periodic soft weight resets to prevent loss of plasticity, annealing $n$-step returns ($n: 10 \to 3$), and dueling/distributional heads \bookcite{reinforement_learning_kevin_murphy}{Sec.~2.5.4.7, pp.~48--49}. For sparse-reward goal-conditioned tasks, \emph{Hindsight Experience Replay (HER)} relabels failed trajectories with achieved terminal states $g' = s_T$ inside the off-policy buffer $\mathcal{D}$ \bookcite{reinforement_learning_kevin_murphy}{Sec.~2.5.5, p.~49}.
\end{crossbook}

\begin{table}[htbp]
\centering
\caption{Per-step computational complexity, memory footprint, and convergence properties ($d$: parameter dimension; $\abs{\Actions}$: action count; $B$: minibatch size).}
\label{tab:td:complexity}
\small
\begin{tabularx}{\textwidth}{@{}l l l l X@{}}
\toprule
\textbf{Method} & \textbf{Time / Step} & \textbf{Memory} & \textbf{Policy} & \textbf{Linear / Nonlinear Stability} \\
\midrule
First-Visit MC & $O(1)$ per step & $O(T)$ episode & On / Off (IS) & Unbiased; converges under coverage \\
$\mathrm{TD}(0)$ / SARSA & $O(d)$ & $O(d)$ & On-policy & Converges w.p.1 (linear, on-policy) \\
$\mathrm{TD}(\lambda)$ & $O(d)$ & $O(d)$ (trace $\vecb{z}$) & On-policy & Converges w.p.1 (linear, on-policy) \\
Q-Learning & $O(\abs{\Actions} d)$ & $O(d)$ & Off-policy & Diverges in general (Deadly Triad) \\
Gradient-TD (TDC) & $O(d)$ ($\approx 2\times$ TD) & $O(2d)$ ($\vecb{w}, \vecb{v}$) & Off-policy & Converges to $\mathrm{PBE}=0$ (linear) \\
Double DQN + PER & $O(B \cdot \text{fwd/bwd})$ & $O(N)$ buffer & Off-policy & Empirically stabilized via $\mathcal{D}, \vecb{\theta}^-$ \\
\bottomrule
\end{tabularx}
\end{table}

\section{Visual Diagnostics \& Geometric Structure}
\label{sec:td:visuals}

Figure~\ref{fig:td:deadly_triad_and_dqn} maps the causal anatomy of the Deadly Triad against the stabilizing mechanisms of modern deep value-based RL, while Figure~\ref{fig:td:bias_variance_nstep} illustrates the empirical bias--variance trade-off across the $n$-step / $\lambda$ bootstrapping spectrum.

\begin{figure}[htbp]
\centering
\begin{tikzpicture}[
  node distance=1.3cm and 1.2cm,
  >=Stealth,
  every node/.style={font=\small}
]
  % Left: Deadly Triad Triangle
  \node[monoWarnBox, text width=2.85cm, align=center] (fa) {\textbf{Function Approx.}\\$\hat{v}(s,\vecb{w})$ / $Q_{\vecb{\theta}}(s,a)$};
  \node[monoWarnBox, below left=1.45cm and -0.35cm of fa, text width=2.65cm, align=center] (boot) {\textbf{Bootstrapping}\\TD / DP targets};
  \node[monoWarnBox, below right=1.45cm and -0.35cm of fa, text width=2.65cm, align=center] (offp) {\textbf{Off-Policy Data}\\$\mu \ne d^\pi$, $\max_{a'}$};

  \draw[monoArrow, monographCrimson, <->] (fa) -- (boot);
  \draw[monoArrow, monographCrimson, <->] (fa) -- (offp);
  \draw[monoArrow, monographCrimson, <->] (boot) -- (offp);

  \node[below=3.15cm of fa, text width=5.4cm, align=center, font=\footnotesize\color{monographCrimson}] (triad_lbl) {\textbf{Deadly Triad $\implies$ Instability}\\Non-contractive $\Pi\Bellman_\pi$ (Baird divergence)};

  % Right: Stabilizers (positioned cleanly to the right of offp)
  \node[monoIndigoBox, right=1.0cm of offp, yshift=0.8cm, text width=4.8cm, align=left] (stab2) {\textbf{2. Frozen Target Net ($\vecb{\theta}^-$) \& Double Q}\\Stabilizes moving targets; decouples $\argmax$ selection from evaluation.};
  \node[monoTealBox, above=0.35cm of stab2, text width=4.8cm, align=left] (stab1) {\textbf{1. Experience Replay ($\mathcal{D}$ / PER)}\\Breaks temporal correlation; prioritizes high-$\abs{\delta_i}$ transitions.};
  \node[monoBox, below=0.35cm of stab2, text width=4.8cm, align=left] (stab3) {\textbf{3. Dueling ($V+A$) \& Gradient-TD}\\Separates state baseline; TDC enforces true SGD on $\mathrm{PBE}(\vecb{w})$.};

  \draw[monoArrow, monographNavy, dashed] (offp.east) -- (stab2.west);
\end{tikzpicture}
\caption{Anatomy of the Deadly Triad (left) and the complementary algorithmic stabilizers deployed in Gradient-TD and Rainbow DQN (right) \bookcite{reinforcement_learning_introduction_richard_s_su}{Ch.~11, PDF p.~286} \bookcite{reinforement_learning_kevin_murphy}{Sec.~2.5.2--2.5.4, pp.~42--48}.}
\label{fig:td:deadly_triad_and_dqn}
\end{figure}

\begin{figure}[htbp]
\centering
\begin{tikzpicture}
\begin{axis}[
  width=0.84\textwidth,
  height=5.6cm,
  xlabel={Step size $\alpha$},
  ylabel={RMS Value Error},
  xmin=0, xmax=1.0,
  ymin=0.22, ymax=0.58,
  grid=major,
  grid style={dashed, gray!30},
  legend style={at={(0.5,1.03)}, anchor=south, legend columns=2, font=\footnotesize, fill=white, /tikz/every even column/.append style={column sep=0.4cm}},
  tick label style={font=\footnotesize},
  label style={font=\small}
]
  % lambda = 0 (TD(0)) - low variance, higher bias at small alpha, stable at moderate alpha
  \addplot[color=monographNavy, very thick, mark=o, mark size=1.5pt] coordinates {
    (0.05, 0.50) (0.1, 0.44) (0.2, 0.37) (0.3, 0.34) (0.4, 0.33) (0.5, 0.34) (0.6, 0.36) (0.8, 0.43) (1.0, 0.53)
  };
  \addlegendentry{$\lambda = 0$ ($\mathrm{TD}(0)$, high bias / low variance)}

  % lambda = 0.8 (Intermediate) - best bias-variance sweet spot
  \addplot[color=monographTeal, very thick, mark=square*, mark size=1.5pt] coordinates {
    (0.05, 0.47) (0.1, 0.38) (0.2, 0.29) (0.3, 0.25) (0.4, 0.24) (0.5, 0.25) (0.6, 0.27) (0.8, 0.34) (1.0, 0.48)
  };
  \addlegendentry{$\lambda = 0.8$ (Optimal intermediate trade-off)}

  % lambda = 0.95
  \addplot[color=monographIndigo, thick, dashed, mark=triangle*, mark size=1.5pt] coordinates {
    (0.05, 0.46) (0.1, 0.37) (0.2, 0.30) (0.3, 0.28) (0.4, 0.29) (0.5, 0.32) (0.6, 0.36) (0.8, 0.47)
  };
  \addlegendentry{$\lambda = 0.95$ (Near-MC)}

  % lambda = 1.0 (MC) - unbiased but high variance, diverges/degrades rapidly with alpha
  \addplot[color=monographCrimson, very thick, dotted, mark=diamond*, mark size=1.8pt] coordinates {
    (0.05, 0.48) (0.1, 0.43) (0.2, 0.40) (0.3, 0.42) (0.4, 0.47) (0.5, 0.54)
  };
  \addlegendentry{$\lambda = 1$ (Monte Carlo, unbiased / high variance)}
\end{axis}
\end{tikzpicture}
\caption{Bias--variance trade-off across $\lambda \in [0, 1]$ on the 19-state random walk benchmark (qualitative profile matching \bookcite{reinforcement_learning_introduction_richard_s_su}{Fig.~12.8, PDF p.~321}). Intermediate $\lambda$ (or $n$-step horizons) consistently achieves the lowest RMS error.}
\label{fig:td:bias_variance_nstep}
\end{figure}

\section{Numerical Pitfalls, Edge Cases \& Failure Modes}
\label{sec:td:pitfalls}

\begin{pitfall}[Unbounded Variance in Ordinary Importance Sampling]
In off-policy MC control, when the behavior policy $b$ explores frequently while $\pi$ is deterministic, the trajectory ratio $\rho_{t:T-1}$ in \eqref{eq:td:is_ratio} can grow exponentially in episode length $T - t$, yielding infinite variance even when the estimator is unbiased \bookcite{reinforcement_learning_introduction_richard_s_su}{Ch.~5, PDF p.~129}. \textbf{Remedy:} Use weighted importance sampling, per-decision IS, or $n$-step Tree Backup / $Q(\sigma)$ methods that avoid IS ratios altogether \bookcite{reinforcement_learning_introduction_richard_s_su}{Ch.~7, PDF p.~174}.
\end{pitfall}

\begin{pitfall}[Naive Residual-Gradient / Bellman Error Minimization]
Replacing semi-gradient TD with direct SGD on the Bellman Error $\mathrm{BE}(\vecb{w}) = \norm{\overline{\delta}_{\vecb{w}}}_\mu^2$ fails in two ways: (1) sampling $s'$ once per step yields a biased gradient of $\E[R + \gamma V(S')]^2$ (the double-sampling problem), and (2) the Bellman Error itself is \emph{not learnable} from observed feature-reward trajectories without knowing the underlying state transition probabilities \bookcite{reinforcement_learning_introduction_richard_s_su}{Ch.~11, PDF pp.~291--298}. \textbf{Remedy:} Minimize the \emph{Projected} Bellman Error via Gradient-TD (TDC/GTD2) or Emphatic-TD.
\end{pitfall}

\begin{pitfall}[Target Network Staleness vs. Loss of Plasticity in Deep Q-Learning]
Updating the frozen target $\vecb{\theta}^-$ too infrequently slows value propagation across long horizons, whereas updating it too frequently reintroduces moving-target oscillations \bookcite{understandingdl_simonprince}{Sec.~19.4.1, p.~401}. Furthermore, high update-to-data (UTD) ratios in sample-efficient DQN can induce dormant neurons and loss of plasticity \bookcite{reinforement_learning_kevin_murphy}{Sec.~2.5.4.7, p.~49}. \textbf{Remedy:} Use Polyak averaging ($\tau \approx 0.005$), Huber/clipped TD losses, layer normalization, and periodic soft parameter resets for high-UTD regimes.
\end{pitfall}

\endgroup
\end{document}
